\section{Introduction}\label{sec:intro}

Audit findings are used twice. They are outputs of completed work, and they
become inputs to the next audit plan. A prior finding can change what the
auditor tests, how much is tested, and which known conditions receive
follow-up. The next finding record is therefore produced partly by what
problems exist and partly by where the audit chose to look.

This paper asks what a record of prior findings can tell us about future
underlying problems when the prior findings themselves change future
observation. The answer depends on which of two jobs the record is being asked
to do. Follow-up asks whether a known condition remains. Discovery asks where
conditions not already represented in the finding history exist. A prior
finding is close to ideal routing information for follow-up: it identifies a
known condition that deserves revisiting, and returning there is what a
risk-based audit is for. It is not sufficient verification that the
condition remains; the follow-up work generates that evidence. Discovery is different. Prior
findings can contain information about discovery, because some entities are
persistently susceptible across the board and one underlying weakness can
affect several kinds of condition. But evidence generated by returning to
known conditions does not by itself identify where unreported conditions lie,
and the finding record alone does not identify how much of that relationship
reflects underlying conditions and how much reflects the observation process
created by the prior findings themselves, because the places without findings
are precisely where absence, lack of observation, and failed detection leave
the same recorded outcome.

The point can be made exact with three quantities. Write $F_t$ for whether a
finding is recorded for a program this year, $C_t$ for whether the underlying
condition actually exists this year, and $f$ for whether a finding was recorded
in the prior period, so $f = 1$ after a finding year and $f = 0$ after a clean
year. A finding is recorded only if the condition exists and the audit detects
it, so the chance of a finding this year, given last year's record, factors
into two pieces:
%
\begin{equation}\label{eq:long}
  P(F_t = 1 \mid F_{t-1} = f)
  \;=\;
  P(C_t = 1 \mid F_{t-1} = f)
  \;\times\;
  P(F_t = 1 \mid C_t = 1,\; F_{t-1} = f) .
\end{equation}
%
Give the three pieces short names. $q_f$ is the observed probability of a
current finding after prior history $f$: what the record shows. $r_f$ is the
probability the underlying condition is actually present after history $f$:
what a planner wants to know. $d_f$ is the probability the audit detects the
condition when it is present, after history $f$: how hard and where the audit
looked. Equation~\eqref{eq:long} then reads
%
\begin{equation}\label{eq:qrd}
  q_f \;=\; r_f \, d_f .
\end{equation}
%
In audit terms: what appears in the finding record depends on two things, how
often the problem is actually there and how often the audit catches it when it
is there. A prior finding can move either one. Remediation changes $r$.
Follow-up changes $d$. The record shows only $q$. Fig.~\ref{fig:diagram}
draws the sequence, together with the feedback arrow that gives a prior finding
its second use: this year's record enters next year's risk assessment,
follow-up plan and scope, and so changes next year's $d$. When that arrow is
active, a difference in $q$ between history groups is a difference in $r$, a
difference in $d$, or both, and the record does not say which. All three
quantities compare groups defined by prior recorded history; they are not
before-and-after effects of remediation or follow-up on a given program, which
are the audit mechanisms that motivate them.

One assumption is maintained throughout: a recorded finding reflects a real
condition. If findings can be recorded where no condition exists, at rate
$a_f$ after history $f$, the recorded rate becomes $q_f = r_f d_f +
(1 - r_f)\, a_f$. Equation~\eqref{eq:qrd} and every bound built on it hold
under no false positives, or under a modeled restriction on $a_f$ carried
through the algebra (Appendix~\ref{app:id}).

Verification bias, imperfect detection, endogenous inspection and selective
labels are established problems
\citep{ransohoff1978,begg1983,mackenzie2002,heckman1978,lakkaraju2017,waymire2018}.
Audit research and practice use prior findings both to measure risk and to
determine subsequent verification. The audit-finding studies most directly related to this paper document
recurrence and predictive association, and recognize that a reported finding
depends on the condition and its detection
\citep{keating2005,petrovits2011,rice2012,waymire2018}, but do not
separately identify the latent-risk and observation components of the
finding history. This paper shows that the two uses make the resulting finding
history endogenous to the audit design, derives what remains identified, and
shows how audit design can recover the missing information. Four
contributions follow.
\emph{Identification:} Section~\ref{sec:worlds} shows exactly why a finding
history cannot, on its own, separate persistence of underlying problems from
persistence of observation, by building three programs with the same record
whose condition persists, is unchanged, or has been halved. \emph{Positive
bounds:} Section~\ref{sec:bounds} shows that if an existing condition is at
least as likely to be detected after a prior finding as after a clean history,
the observed finding risk ratio is an upper bound on the underlying-condition
risk ratio, and that a defensible ceiling on the follow-up detection
multiplier makes the bound two-sided; Section~\ref{sec:operational} sets out
how an organization can learn or bound that multiplier. \emph{Measurement:}
Section~\ref{sec:fac} distinguishes repeat-only, new-only and mixed outcomes
and shows why a recurrence-share statistic can range from \nZetaLo{} to
\nZetaHi{} percent without any change in the data. \emph{Empirical
illustration:} in Federal Audit Clearinghouse data, prior findings predict
both recurrence and materially more findings the auditor did not classify as
repeats, among programs audited as major programs in both adjacent years,
while the record itself cannot identify why.

The Single Audit is the illustration because history-dependent observation is
written into its rules rather than inferred from auditor behavior;
Section~\ref{sec:regime} sets those rules out before the data are used. The
argument itself is general: regulatory enforcement, quality assurance,
clinical surveillance and tax examination all produce records with the same
property. Proofs are in Appendix~\ref{app:id}.

\begin{figure}[tbp]
\centering
\includegraphics[width=\textwidth]{ex1_diagram.pdf}
\caption{\textbf{A finding history is also an observation history.} A recorded
finding ($F_t = 1$) requires the condition to exist ($C_t = 1$), the audit to
examine the area where it is, and the auditor to recognize it there. The
detection probability $d$ covers both steps together, given that the
condition exists: selection into observation as well as recognition.
The absence of a recorded finding is therefore three different things wearing
the same clothes: the condition was not there, or it was never looked for, or
it was looked for and missed. The two feedback arrows are what make the
record hard to read. The prior recorded finding $F_{t-1}$ feeds back two
ways: through corrective action, if any is taken, into next period's
condition $C_t$, and
through next period's risk assessment, follow-up plan and scope into next
period's detection probability $d$. Equation~\eqref{eq:long} is this diagram
written as a product, with $F_{t-1}$ inside both factors. In the Single Audit the feedback is written into the regulation rather
than inferred from behavior (Section~\ref{sec:regime}).}
\label{fig:diagram}
\end{figure}
